\documentclass[10pt,a4paper]{amsart}
\usepackage[margin=19mm]{geometry}
\usepackage{amsmath,amssymb,mathtools}
\usepackage[scr=rsfs,cal=euler]{mathalfa}
\usepackage{xfrac}
\usepackage[unicode,hidelinks,bookmarks=false]{hyperref}
\newcommand{\ie}{i.e.\ }
\newcommand{\cf}{cf.\ }
\newcommand{\resp}{respectively\ }
\newcommand{\Subsection}{\subsection}
\providecommand{\nobreakdash}{\nobreak\hbox{-}}
\setlength{\emergencystretch}{2em}
\multlinegap0pt
\usepackage{amssymb}
\usepackage{enumerate}
\usepackage{bm}
\usepackage{xcolor}
\newtheorem{prop}{Proposition}
\newtheorem{thm}{Theorem}
\newtheorem{lem}{Lemma}
\DeclareMathOperator{\Hom}{Hom}
\DeclareMathOperator{\re}{Re}
\title{Notes on acceptable bundles I}
\author{Osamu Fujino, Taro Fujisawa,, Takashi Ono}
\date{Source: arXiv:2511.00760v2 (CC BY 4.0)}
\begin{document}
\maketitle

\noindent\emph{Source and licence.} Notes on acceptable bundles I. Version arXiv:2511.00760v2 (CC BY 4.0), distributed by arXiv under the Creative Commons Attribution 4.0 International licence, http://creativecommons.org/licenses/by/4.0/, as recorded in the arXiv licence metadata for this exact version and shown on the abstract page https://arxiv.org/abs/2511.00760v2. Full licence notice: \texttt{licenses/arxiv-2511.00760-CC-BY-4.0.txt}.\par\smallskip

\noindent\textit{Editorial scope.} Counted unit: the article's thm p-thm7.5 (source0.tex lines 1976--1984). The article's complete original proof of it is delivered with it (source0.tex lines 1990--2178, one \texttt{proof} environment of 189 lines, which the article places away from the statement). No mathematics is written, repaired or re-derived: the statement and the proof are the article's own lines, in the article's own notation, and no retained line is substituted or re-flowed.  Retained uncounted as the source-local context the proof needs: lines 233--246; lines 367--375; lines 623--632; lines 858--879; lines 947--955; lines 1004--1010; lines 1925--1929.  Nothing was omitted from any retained span, so the package carries no omission record.  No retained line contains a citation, so the wrapper carries no bibliography and no bibliography entry.  No retained line contains a figure environment, an image-inclusion command or drawing code, so no image is delivered and the package declares no asset dependency.  Editorial formatting only, in addition: an AMS \texttt{amsart} preamble that restates the article's own declarations and macros used by the retained lines, namely the environments \texttt{prop}, \texttt{thm}, \texttt{lem} on separate counters, together with the standard packages those lines need.  The retained environments print in this document as Proposition 1, Theorem 1, Lemma 1 in order of appearance, whereas the article prints the same items with the numbering of its own full document.  The complete native source of the article as distributed in the arXiv e-print for this exact version is bundled unchanged as \texttt{sources/188274/source0.tex}.
\begin{prop}[Proposition \ref{p-prop4.1}]\label{p-prop1.5} 
Let $L$ be a holomorphic line bundle 
on $\Delta^*$ and let $h$ be a Hermitian metric on $L$ such that
\[
-C \cdot \omega_P \leq \sqrt{-1} \Theta_h(L) \leq C \cdot \omega_P
\]
holds on $\Delta^*$ for some constant $C > 0$.
That is, writing
\[
\sqrt{-1} \Theta_h(L) = f(z) \cdot \omega_P,
\]
we have $|f(z)| \leq C$ on $\Delta^*$. 
Then ${}_a L$ is a holomorphic line bundle for every $a \in \mathbb{R}$.
\end{prop}

\begin{thm}[Determinant bundles, see Theorem \ref{p-thm7.5}]\label{p-thm1.10}
Let $(E, h)$ be an acceptable vector bundle on $\Delta^*$. Then 
the determinant bundle $(\det E, \det h)$ is an acceptable 
line bundle on $\Delta^*$, and
\[
\det({}_a E) = {}_{\gamma({}_a E)} \det E
\]
holds for every $a \in \mathbb{R}$.
\end{thm}

\begin{lem}\label{p-lem2.2}
Let $(E, h)$ be an acceptable vector bundle on $\Delta^*$. 
Then the dual bundle $(E^\vee, h^\vee)$ and the 
determinant line bundle $(\det E, \det h)$ are also acceptable. 

Let $(E_1, h_1)$ and $(E_2, h_2)$ be acceptable 
vector bundles on $\Delta^*$. 
Then the tensor product $(E_1 \otimes E_2, h_1 \otimes h_2)$ 
and the Hom bundle $(\Hom(E_1, E_2), h_1^\vee \otimes h_2)$ are acceptable. 
\end{lem}

\begin{lem}[Lelong number]\label{p-lem3.1}
Let $u$ be a subharmonic function on $\Delta$. Then we have
\begin{equation}\label{p-eq3.1} 
\lim_{r \to 0} \int_{\Delta(0, r)} 
\frac{\sqrt{-1}}{\pi} \partial \overline{\partial} u 
= \liminf_{z \to 0} \frac{u(z)}{\log |z|}.
\end{equation}
We define
\[
\nu(u, 0) := \liminf_{z \to 0} \frac{u(z)}{\log |z|}
\]
and call it the {\em Lelong number} of $u$ at $0$. 
Note that the expression $\partial \overline{\partial} u$ is 
understood in the sense of currents.

Let $u_1$ and $u_2$ be subharmonic 
functions on $\Delta$. Then $u_1 + u_2$ is also subharmonic on $\Delta$.  
By \eqref{p-eq3.1}, we have the identity
\[
\nu(u_1 + u_2, 0) = \nu(u_1, 0) + \nu(u_2, 0).
\]
\end{lem}

\begin{lem}\label{p-lem3.3}
Let $g$ be a holomorphic function on $\Delta^*$. Assume that
\[
\re g(z) \leq C (-\log |z|)
\]
holds on $\Delta^*$ for some constant 
$C > 0$. Then $g$ extends holomorphically to 
the origin; that is, the origin is a removable singularity of $g$.
\end{lem}

\begin{prop}[Proposition \ref{p-prop1.5}]\label{p-prop4.1} 
Let $(L, h)$ be an acceptable line bundle 
on $\Delta^*$. 
Then ${}_\alpha L$ 
is a holomorphic line bundle 
on $\Delta$ for every $\alpha \in \mathbb R$. 
\end{prop}

\begin{lem}\label{p-lem7.3}
\[
-\infty<\liminf _{z\to 0} \frac{\log \det H(\bm v)}{\log |z|}<\infty. 
\]
\end{lem}

\begin{thm}[Determinant bundles, see Theorem \ref{p-thm1.10}]\label{p-thm7.5}
Let $(E, h)$ be an acceptable vector bundle on $\Delta^*$. 
Let $\alpha$ be any real number. 
Then 
\[
\det ({}_\alpha E)={}_{\gamma ({}_\alpha E)} \det E
\] 
holds. 
\end{thm}

\begin{proof}[Proof of Theorem \ref{p-thm7.5}]
Let $\bm v:=\{v_1, \ldots, v_r\}$ be a frame of ${}_\alpha E$ on $\Delta$. 
We put 
\[
H(h, \bm v)(z):=\left(h(v_i, v_j)(z)\right)_{i, j}. 
\] 
We note that 
\[
\det ({}_\alpha E)=\mathcal O_\Delta v_1\wedge \cdots \wedge v_r
\] 
and 
\begin{equation}\label{p-eq7.2}
\det E=\mathcal O_{\Delta^*}v_1\wedge \cdots \wedge v_r\simeq 
\mathcal O_{\Delta^*}. 
\end{equation} 
We consider $L:=\det E$. Let $h_L$ be the induced metric on $L$. 
Note that $(L, h_L)$ is an acceptable line bundle on $\Delta^*$ 
by Lemma \ref{p-lem2.2} since $(E, h)$ is acceptable. 
By using the trivialization \eqref{p-eq7.2}, 
we argue as in the proof of Proposition \ref{p-prop4.1}. 
In this setting, 
\[
h_L=|\cdot |^2 e^{-2\varphi_\alpha}
\] 
with 
\begin{equation}\label{p-eq7.3}
\begin{split}
e^{-2\varphi_\alpha} &=h_L(v_1\wedge \cdots \wedge v_r, v_1\wedge \cdots \wedge v_r) 
\\&=\det H(\bm v). 
\end{split}
\end{equation} 
Thus 
\[
\varphi_\alpha=-\frac{1}{2} \log \det H(\bm v). 
\] 
Therefore, 
\begin{equation}\label{p-eq7.4}
\liminf_{z\to 0} \frac{-\varphi_\alpha(z)}{\log |z|} 
=\frac{1}{2} \liminf_{z\to 0} \frac{\log \det H(\bm v)}{\log |z|} 
=-\gamma ({}_\alpha E). 
\end{equation} 
Note that $(L, h_L)$ is an acceptable line bundle on $\Delta^*$. 
As in the proof of Proposition \ref{p-prop4.1}, we can write 
\begin{equation}\label{p-eq7.5}
2\varphi_\alpha +\chi (N)=2\psi_1+c_1\log |z|^2+2\re g_1(z)
\end{equation}
and 
\begin{equation}\label{p-eq7.6}
-2\varphi_\alpha +\chi (N)=2\psi_2+c_2\log |z|^2+2\re g_2(z), 
\end{equation} 
where $\psi_1$ and $\psi_2$ are subharmonic functions on $\Delta$ and 
$g_1$ and $g_2$ are holomorphic functions on $\Delta^*$. 
By \eqref{p-eq7.5}, we have 
\[
-\re g_1(z)=\psi_1+c_1\log |z|-\varphi_\alpha -\chi (N). 
\]
Hence we have 
\[
\frac{-\re g_1(z)}{\log |z|}=\frac{\psi_1(z)}{\log |z|}
+c_1+\frac{-\varphi_\alpha(z)}{\log |z|}
+\frac{-\chi(N)}{\log|z|}. 
\] 
Therefore, we obtain 
\[
\begin{split}
\liminf_{z\to 0}\frac{-\re g_1(z)}{\log |z|}&\geq 
\liminf_{z\to 0}\frac{\psi_1(z)}{\log |z|}
+c_1+\liminf_{z\to 0}\frac{-\varphi_\alpha(z)}{\log|z|}
+\liminf_{z\to 0}\frac{-\chi(N)}{\log|z|}
\\ &=\nu(\psi_1, 0)+c_1-\gamma({}_\alpha E) \\ 
&>-\infty. 
\end{split}
\] 
Here we used Lemma \ref{p-lem7.3}. 
Thus there exists some $C>0$ such that 
\[
\frac{-\re g_1(z)}{\log |z|}\geq -C
\] 
holds over some open neighborhood of $0$. 
This implies that 
\[
\re\left(-g_1(z)\right)\leq C\left(-\log |z|\right)
\] 
holds around $0$. 
By Lemma \ref{p-lem3.3}, we see that 
$g_1$ is holomorphic on $\Delta$. 
Hence we have 
\[\nu_1+c_1-\gamma ({}_\alpha E)\leq 0=\liminf_{z\to 0} \frac{-\re g_1(z)}{\log |z|}, 
\] 
where $\nu_1:=\nu(\psi_1, 0)$. 
Note that $e^{g_1(z)}$ is a nowhere vanishing holomorphic 
function on $\Delta$. 
By replacing $v_1$ with $e^{g_1(z)}v_1$,  
we may assume that
\begin{equation}\label{p-eq7.7}
2\varphi_\alpha + \chi(N) = 2\psi_1 + c_1 \log |z|^2
\end{equation}
and
\begin{equation}\label{p-eq7.8}
-2\varphi_\alpha + \chi(N) = 2\psi_2 + c_2 \log |z|^2 + 2\re g_2(z)
\end{equation}
hold, after replacing $g_2$ accordingly. 
By \eqref{p-eq7.7} and \eqref{p-eq7.8}, we have 
\[
2\chi (N)=2\psi_1+2\psi_2+(c_1+c_2)\log |z|^2+2\re g_2(z) 
\] 
holds. 
Note that $\chi(N)$, $\psi_1$, $\psi_2$, and $\log |z|^2$ are subharmonic 
functions on $\Delta$. 
We have 
\[
\frac{-\re g_2(z)}{\log |z|}=\frac{\psi_1(z)}{\log |z|}
+\frac{\psi_2(z)}{\log |z|}+c_1+c_2
+\frac{-\chi(N)}{\log|z|}. 
\] 
Therefore, we obtain 
\[
\begin{split}
\liminf_{z\to 0}\frac{-\re g_2(z)}{\log |z|}&\geq 
\liminf_{z\to 0}\frac{\psi_1(z)}{\log |z|}
+\liminf_{z\to 0}\frac{\psi_2(z)}{\log |z|}
+c_1+c_2
+\liminf_{z\to 0}\frac{-\chi(N)}{\log|z|}
\\ &=\nu(\psi_1, 0)+\nu(\psi_2, 0)+c_1+c_2. 
\end{split}
\] 
Thus there exists some $C>0$ such that 
\[
\frac{-\re g_2(z)}{\log |z|}\geq -C
\] 
holds over some open neighborhood of $0$. 
This implies that 
\[
\re\left(-g_2(z)\right)\leq C\left(-\log |z|\right)
\] 
holds around $0$. 
By Lemma \ref{p-lem3.3}, we see that 
$g_2$ is holomorphic on $\Delta$. 
In particular, $\re g_2(z)$ is harmonic on $\Delta$. 
By replacing $\psi_2$ with $\psi_2 +\re g_2(z)$, we can finally assume that 
\begin{equation}\label{p-eq7.9}
2\varphi_\alpha +\chi (N)=2\psi_1+c_1\log |z|^2
\end{equation}
and 
\begin{equation}\label{p-eq7.10}
-2\varphi_\alpha +\chi (N)=2\psi_2+c_2\log |z|^2
\end{equation} 
hold. 
By \eqref{p-eq7.10}, we obtain 
\[
-\gamma ({}_\alpha E)=\nu_2+c_2,  
\] 
where $\nu_2:=\nu(\psi_2, 0)$. 
By \eqref{p-eq7.9} and \eqref{p-eq7.10}, we have 
\[
\chi (N)=\psi_1+\psi_2+(c_1+c_2)\log |z|. 
\] 
Therefore, by Lemma \ref{p-lem3.1}, we obtain 
\[
0=\nu_1+\nu_2+c_1+c_2. 
\] 
This means that 
\[\nu _1+c_1=-(\nu_2+c_2)=\gamma ({}_\alpha E). 
\]
By \eqref{p-eq7.9}, 
we have 
\begin{equation}\label{p-eq7.11}
\liminf_{z\to 0} \frac{\varphi_\alpha(z)}{\log |z|}=\nu_1+c_1 =
\gamma ({}_\alpha E).   
\end{equation} 
Thus, by \eqref{p-eq7.4} and \eqref{p-eq7.11}, we get 
\[
\lim_{z\to 0} \frac{\varphi_\alpha(z)}{\log |z|}=
\gamma ({}_\alpha E), \quad \lim _{z\to 0} \frac{\psi_1(z)}{\log |z|}=\nu_1, 
\quad \text{and}\quad 
\lim _{z\to 0} \frac{\psi_2(z)}{\log |z|}=\nu_2.  
\] 
As in the proof of Proposition \ref{p-prop4.1}, we obtain 
\begin{equation}\label{p-eq7.12}
{}_\beta L={}_\beta (\det E)=\mathcal O_{\Delta} \cdot 
z^{-\lfloor \beta -\gamma ({}_\alpha E)\rfloor}v_1\wedge \cdots 
\wedge v_r. 
\end{equation} 
In particular, 
\[
\det ({}_\alpha E)={}_{\gamma ({}_\alpha E)}\det E. 
\] 
We finish the proof of Theorem \ref{p-thm7.5}. 
\end{proof}

\end{document}
